\section{Task, Ontology, and Benchmark}
\label{sec:benchmark}

\subsection{The TITAN Ontology}
\label{sec:benchmark:ontology}

The TITAN knowledge graph is derived from MITRE ATT\&CK CITAZIONE and organizes CTI entities into nine types spanning three conceptual roles: an \emph{adversarial} side --- Attack Pattern, Malware, Tool, Intrusion Set, and Campaign --- capturing techniques, software, and the actors who use them; a \emph{defense and observation} side --- Course of Action, Data Component, and Data Source --- capturing mitigations and detection telemetry; and a single \emph{target} type, Asset, capturing the systems techniques act on. Every relation in the graph is \emph{typed}, so that an edge label alone disambiguates its target type even when several relations share the same verb (e.g.\ \texttt{uses\_attack\_pattern} vs.\ \texttt{uses\_malware}), and almost every relation is \emph{bidirectional}, with a relation and its inverse both present as first-class edges of consistent, complementary semantics ---
\emph{malware} $\xrightarrow{\texttt{uses\_attack\_pattern}}$ \emph{attack pattern} and \emph{attack pattern} $\xrightarrow{\texttt{used\_by\_malware}}$ \emph{malware} are both present --- which is what lets the path planner begin reasoning from any node in the graph, not only from a canonical "subject" position. The one exception is \texttt{targets} (attack pattern $\to$ asset), kept one-directional since an asset is a terminal object of reasoning rather than a waypoint one would reason \emph{from}. Figure~\ref{fig:ontology} shows the type inventory grouped by role together with a representative set of typed relations; full per-type node counts and the complete relation inventory are given in Appendix~\ref{sec:appendix} (Table~\ref{tab:ontology_counts}).

\begin{figure*}[t]
\centering
\begin{tikzpicture}[
  adv/.style={draw=titanAdversarial, rounded corners=3pt, align=center,
              minimum height=6mm, minimum width=15mm, inner sep=1.5pt,
              font=\scriptsize\bfseries, line width=0.6pt,
              fill=titanAdversarial!10, text=titanAdversarial!80!black},
  defbox/.style={draw=titanDefense, rounded corners=3pt, align=center,
              minimum height=6mm, minimum width=15mm, inner sep=1.5pt,
              font=\scriptsize\bfseries, line width=0.6pt,
              fill=titanDefense!10, text=titanDefense!80!black},
  tgt/.style={draw=titanNeutral, rounded corners=3pt, align=center,
              minimum height=6mm, minimum width=15mm, inner sep=1.5pt,
              font=\scriptsize\bfseries, line width=0.6pt,
              fill=titanNeutralLight!50, text=titanNeutral},
  fwd/.style={-{Latex[length=1.1mm]}, draw=titanNeutral, thin},
  badge/.style={circle, draw=titanNeutral, fill=white, minimum size=3.2mm,
                inner sep=0pt, font=\tiny\bfseries, text=titanNeutral},
  swatch/.style={rounded corners=1pt, minimum width=2.4mm, minimum height=2.4mm, inner sep=0pt},
  legtxt/.style={font=\tiny, text=titanNeutral, anchor=west},
]
% --- adversarial cluster ---
\node[adv] (is)   at (0,1.4)    {Intrusion Set};
\node[adv] (cam)  at (0,-1.4)   {Campaign};
\node[adv] (mal)  at (2.6,2.0)  {Malware};
\node[adv] (tool) at (2.6,-0.3) {Tool};
\node[adv] (ap)   at (5.3,0.75) {Attack Pattern};

% --- defense / observation cluster ---
\node[defbox] (coa) at (8.0,2.0)  {Course of\\Action};
\node[defbox] (dc)  at (8.0,0.3)  {Data\\Component};

% --- target ---
\node[tgt] (asset) at (5.3,-1.8) {Asset};

% --- typed, bidirectional edges: each pair is two single-headed arcs (one
% per direction), each carrying a letter badge for its relation (legend
% below). uses/used\_by cover several source types (Section~\ref{sec:benchmark:ontology});
% targets is the one relation with no inverse.
\draw[fwd, bend left=12] (is)   to node[badge, pos=0.5] {A} (mal);
\draw[fwd, bend left=12] (mal)  to node[badge, pos=0.5] {B} (is);

\draw[fwd, bend left=12] (cam)  to node[badge, pos=0.5] {I} (is);
\draw[fwd, bend left=12] (is)   to node[badge, pos=0.5] {E} (cam);

\draw[fwd, bend left=12] (mal)  to node[badge, pos=0.5] {A} (ap);
\draw[fwd, bend left=12] (ap)   to node[badge, pos=0.5] {B} (mal);

\draw[fwd, bend left=12] (tool) to node[badge, pos=0.5] {A} (ap);
\draw[fwd, bend left=12] (ap)   to node[badge, pos=0.5] {B} (tool);

\draw[fwd, bend left=12] (coa)  to node[badge, pos=0.5] {F} (ap);
\draw[fwd, bend left=12] (ap)   to node[badge, pos=0.5] {G} (coa);

\draw[fwd, bend left=12] (dc)   to node[badge, pos=0.5] {C} (ap);
\draw[fwd, bend left=12] (ap)   to node[badge, pos=0.5] {D} (dc);

\draw[fwd] (ap) to node[badge, pos=0.5] {H} (asset);

% --- legend: role swatches, then the letter code (3x3) ---
\node[swatch, fill=titanAdversarial!10, draw=titanAdversarial] (legadv) at (0,-2.35) {};
\node[legtxt, right=0.8mm of legadv] {Adversarial};
\node[swatch, fill=titanDefense!10, draw=titanDefense] (legdef) at (2.9,-2.35) {};
\node[legtxt, right=0.8mm of legdef] {Defense \& Obs.};

\node[badge] (lA) at (0,-2.8)   {A}; \node[legtxt, right=0.8mm of lA] {uses};
\node[badge] (lB) at (2.9,-2.8) {B}; \node[legtxt, right=0.8mm of lB] {used by};
\node[badge] (lC) at (5.8,-2.8) {C}; \node[legtxt, right=0.8mm of lC] {detects};
\node[badge] (lD) at (0,-3.2)   {D}; \node[legtxt, right=0.8mm of lD] {detected by};
\node[badge] (lE) at (2.9,-3.2) {E}; \node[legtxt, right=0.8mm of lE] {is responsible};
\node[badge] (lF) at (5.8,-3.2) {F}; \node[legtxt, right=0.8mm of lF] {mitigates};
\node[badge] (lG) at (0,-3.6)   {G}; \node[legtxt, right=0.8mm of lG] {mitigated by};
\node[badge] (lH) at (2.9,-3.6) {H}; \node[legtxt, right=0.8mm of lH] {targets};
\node[badge] (lI) at (5.8,-3.6) {I}; \node[legtxt, right=0.8mm of lI] {attributed to};
\end{tikzpicture}
\caption{TITAN Ontology, by role. Each labeled arc is one typed relation (legend below); \texttt{uses}/\texttt{used~by} letters cover several source types (e.g.\ malware, tool $\to$ attack pattern), spelled out in full in Appendix~\ref{sec:appendix}, which also lists Data Source and the remaining relation instances omitted here for space.}
\label{fig:ontology}
\end{figure*}

\subsection{Dataset Construction}
\label{sec:benchmark:dataset}

The benchmark is generated from 454 hand-written question templates, each carrying typed placeholders (e.g.\ \texttt{[malware]}) and a gold relational path over the ontology; because a path depends only on the node \emph{types} it touches, not on the specific entities filling them, one template yields many training instances as its placeholders are bound to different graph entities. Templates are paraphrased with an LLM (CITAZIONE) to add linguistic variety and discourage the planner from keying on surface template patterns rather than the underlying relation. Some questions have no single explicit starting node --- e.g.\ \emph{"List all ransomware detection strategies"}, where every ransomware instance is a valid start --- and for these the gold path uses one of the four operators (\texttt{filter}, \texttt{select}, \texttt{exec\_common}, \texttt{exec\_difference}) introduced in Section~\ref{sec:methodology:dsl}. Each instance also carries a chain-of-thought rationale that walks the path one relation at a time before stating the final program, used for the CoT reasoning-style condition (Section~\ref{sec:methodology:regimes}).

\subsection{Splits and Statistics}
\label{sec:benchmark:splits}

The benchmark is split into train (65{,}529 rows), validation (3{,}855), and test (15{,}627), for a total of 85{,}011 question--path pairs. The test split is template- and skeleton-disjoint from train and validation --- no test question shares its underlying template, nor a train row's relation skeleton, with anything the model is trained on --- which we verify independently by re-deriving each row's template and skeleton keys directly from its path text rather than trusting a stored split label. Table~\ref{tab:dataset_profile} breaks each split down by path length and by operator; the test split is, by construction, weighted toward longer paths and rarer operators than train (e.g.\ \texttt{select} and \texttt{exec\_difference} each make up roughly twice the share of test that they do of train), which makes it a meaningfully harder distribution, not merely a held-out sample of the same one.

\begin{table*}[t]
\centering
\small
\setlength{\tabcolsep}{4pt}
\begin{tabular}{lrrrr|rrrrr}
\toprule
& \multicolumn{4}{c}{\textit{Path length}} & \multicolumn{5}{c}{\textit{Operator}} \\
\cmidrule(lr){2-5}\cmidrule(lr){6-10}
\textbf{Split} & L1 & L2 & L3 & L4+ & filter & select & com. & diff. & plain \\
\midrule
Train & 27{,}048 & 26{,}247 & 6{,}535 & 5{,}699 & 18{,}634 & 951 & 1{,}336 & 176 & 44{,}432 \\
Val   & 1{,}619  & 1{,}490  & 411    & 335    & 1{,}091  & 56  & 80    & 9   & 2{,}619 \\
Test  & 5{,}984  & 5{,}425  & 1{,}732 & 2{,}486 & 4{,}599  & 485 & 44    & 150 & 10{,}349 \\
\bottomrule
\end{tabular}
\caption{Dataset profile by split (\texttt{com.}~=~\texttt{exec\_common}, \texttt{diff.}~=~\texttt{exec\_difference}, \texttt{plain}~=~no operator).}
\label{tab:dataset_profile}
\end{table*}
